\documentclass[12pt]{article}
\usepackage{graphicx}
\usepackage{amsmath}
\usepackage{url}
\usepackage{booktabs}
\usepackage[margin=1in]{geometry}
\usepackage{natbib}
\usepackage{xurl}
\usepackage{setspace}
\doublespacing

\title{Uneven Legalization, Spatial Demand Spillovers, and the Regulation of Entry:\\
Evidence from U.S.\ Cannabis Markets}
\author{Mahyar Rezazadeh \\ Babak Sahragard}
\date{August 2026}

\begin{document}
\maketitle

\begin{singlespace}
\begin{abstract}
\noindent Uneven state-level legalization of recreational cannabis creates a spatial demand spillover: counties
near a still-prohibited neighboring state face captive cross-border demand that interior counties of the same
state do not. We ask two questions about this spillover. First, does it produce a border premium in legal
cannabis markets. Second, does entry regulation determine which margin absorbs that premium -- existing
sellers' revenue, or the number of sellers and the market's total size. Using a county-year panel spanning
Colorado, Washington, Oregon, Montana, and New Mexico (2014--2026), with a border classification that is
time-varying rather than a fixed snapshot -- accounting for New Mexico's 2022 legalization, Oregon's 2015
legalization, and Oregon's own 2017--2018 transition relative to California and Nevada -- we find a robust
border premium on revenue per dispensary (+48\%, county-clustered $p=0.11$), revenue per capita (+193\%,
$p<0.001$), and dispensary density (+101\%, $p<0.001$). Whether entry is capped, using an institutional
classification that treats binding local-government control over licensing as equivalent to a state cap, does
not moderate the revenue-per-store margin, but it sharply moderates the other two: capped entry reduces the
border premium's multiplier on revenue per capita by 54\% in relative terms (from roughly 193\% to roughly
34\% in level terms, $p=0.009$) and on dispensary density by 56\% in relative terms ($p<0.001$). Oregon's own
within-state entry freeze in January 2022 produces the same-signed, if
underpowered, pattern. A dispensary-level panel built for New Mexico -- the one sample state where individual
firm revenue is observable historically -- confirms the border premium is not a county-aggregation artifact
(+141\% at the level of the individual dispensary, $p=0.001$, significant in every sample year). The results
support a specific mechanism -- entry regulation determines whether spatial
demand spillovers are absorbed by incumbent revenue or by market expansion -- while also reporting, without
forcing a cleaner story, that the mechanism does not extend to the intensive margin the way it was originally
hypothesized to, and that an alternative cut isolating specifically local-government-imposed restrictions
preserves the density result under county clustering but not the other two margins.

\vspace{0.5em}
\noindent\textbf{JEL codes:} L51, H71, R12, L81
\vspace{0.3em}

\noindent\textbf{Keywords:} cannabis, cross-border shopping, market entry, license caps, entry regulation
\end{abstract}
\end{singlespace}

\section{Introduction}
\label{sec:intro}

By the end of 2025, 24 U.S.\ states had legalized recreational cannabis, though Virginia had not yet
established an operating retail market, leaving 23 states with a functioning legal market and 27 states,
plus the remaining territory, where recreational retail sale remained prohibited. This patchwork did not
arrive all at once, and it has not stopped changing: a state that borders only prohibition-states in one year
can border a freshly legal neighbor the next. That timing creates a natural, if underused, source of
identifying variation. A county sitting near a still-prohibited state faces a captive pool of residents who
cannot buy cannabis legally at home; a county in the interior of the same legal state does not. If cannabis
demand is at all price- or access-sensitive across a state line -- which the existing literature on cannabis,
alcohol, and other regulated goods says it is -- that captive demand should show up somewhere in the legal
market's own data.

Where it shows up, however, is not obvious, and that is the question this paper is really about. A dispensary
facing extra demand from across the border can respond in at least two ways: it can sell more at roughly the
same price and margin, raising its own revenue, or new dispensaries can open nearby to meet the extra demand,
raising the county's total revenue and store count without necessarily raising revenue at any one store. Which
of these happens is not a fact about demand; it is a fact about supply-side regulation. A state that caps the
number of licenses, whether by an explicit numerical ceiling or by letting local governments block entry
outright, closes off the second channel by construction. A state that does not cap entry leaves it open. This
paper's central hypothesis is that cross-border demand creates a revenue premium in border counties, and that
entry regulation determines whether that premium is absorbed by existing sellers (revenue per store) or by
market expansion (revenue per capita and dispensary count) -- a specific, testable version of a standard
regulatory-economics prediction: a binding entry restriction forces a demand shock into incumbent rents rather
than into quantity.

We test this using a county-year panel covering five states -- Colorado, Washington, Oregon, Montana, and New
Mexico -- assembled from each state's own revenue-reporting agency (2014--2026, depending on the state's own
retail-launch date), together with an institutional classification of each state's entry-licensing regime built
from the actual statutes and regulations, not from a state's own label for itself. Two features of the design
depart from prior work in this space, including our own earlier work on Colorado and Washington alone. First,
a county's exposure to a nonlegal-state border is treated as time-varying: because several of the sample
states' neighbors themselves legalized cannabis partway through the sample window (New Mexico in 2022, Oregon
in 2015, and, more narrowly, California and Nevada relative to Oregon's own 2017 observations), a border
classification fixed at a single snapshot date misclassifies a meaningful share of county-years. We recompute
it year by year. Second, entry regulation is classified by where the binding constraint on licensing actually
sits, not by whether a state imposes an explicit numerical cap. Colorado has no state-level cap, but roughly
half its counties ban cannabis retail outright through local option, a restriction documented independently by
\citet{PayanBrownSong2021}; Montana has no numerical cap either, but its 2021 enabling statute requires
counties that voted against the state's 2020 legalization initiative to affirmatively opt in before a
dispensary can locate there. Both are coded as capped-entry states for this reason, a coding choice we explain
and defend in Section~\ref{sec:institutional} rather than assert.

We find, first, that the border premium is real and does not depend on regulatory regime: revenue per store,
revenue per capita, and dispensary density are all significantly higher in border counties, pooling across the
five states. Second, entry regulation moderates two of these three margins in exactly the direction the
hypothesis predicts, and does so with more precision than the underlying border effect itself: the interaction
between border exposure and capped entry is a 54\% relative reduction in the border premium's multiplier on
revenue per capita ($p=0.009$) and a 56\% relative reduction on dispensary density ($p<0.001$) -- almost
identical magnitudes, consistent with the
premium moving through the quantity margin specifically, since revenue per capita is arithmetically close to
the product of density and revenue per store. Third, and reported honestly rather than smoothed over, the same
interaction on revenue per store is a precise, stable null: the border premium on revenue per store is
essentially identical in capped and free-entry states. The original hypothesis predicted regulation would push
the premium into the per-store margin specifically when entry is capped; the data show only that regulation
blocks the quantity margin, not that it redirects the premium into price.

The remainder of the paper is organized as follows. Section~\ref{sec:lit} reviews the two literatures this
paper connects. Section~\ref{sec:institutional} details the five states' entry-licensing institutions.
Section~\ref{sec:data} describes the data and the time-varying border construction. Section~\ref{sec:methods}
presents the empirical strategy. Section~\ref{sec:results} reports results. Section~\ref{sec:discussion}
discusses mechanism and limitations, and Section~\ref{sec:conclusion} concludes.

\section{Literature Review}
\label{sec:lit}

Two literatures bear directly on this paper's question, and, to our knowledge, neither has been connected to
the other.

The first documents that cannabis demand crosses state borders when legal status differs. \citet{HaoCowan2020}
show that recreational legalization draws substantial cross-border sales from residents of neighboring
nonlegal states. \citet{HansenMillerWeber2020} show the mechanism operating in the other direction: Washington
retailers near the Oregon border lost roughly 36 percent of their sales when Oregon's own retail market opened
in October 2015, evidence that the earlier premium those counties enjoyed was genuinely a captive-demand
effect. \citet{BaggioChongKwon2020} document a related cross-border effect between cannabis and a different
regulated good, showing that counties in states that legalized medical marijuana reduced alcohol purchases
relative to nearby counties across the state line, using a border-county-pair design in the tradition of
\citet{DubeLesterReich2010}. Outside cannabis, cross-border shopping in response to tax and regulatory
differences is a mature empirical literature; \citet{FribergEtAl2024} is a recent addition, using Norway's
COVID-19 border closures as a natural experiment to show that cross-border shopping measurably erodes domestic
commodity tax revenue.

The second literature studies how entry regulation shapes market structure and rents once a cannabis market
exists, independent of any cross-border demand story. \citet{HollenbeckUetake2021} document that Washington's
capped licensing regime supports meaningful markups. \citet{Thomas2019LicenseQuotas} builds a structural model
of Washington's own license quota and estimates that free entry would raise total surplus by roughly 18
percent relative to the actual capped and partly randomized allocation. This is, as far as we can establish,
the closest existing paper to the entry-regulation half of our own question -- but with no demand shock to
test it against. The general industrial-organization literature on entry supplies the theoretical scaffolding:
\citet{BresnahanReiss1990,BresnahanReiss1991} establish that competitive conduct changes most sharply with the
arrival of a market's second or third entrant; \citet{Berry1992}, \citet{Mazzeo2002}, and \citet{Seim2006}
extend this to settings with firm heterogeneity; \citet{ToivanenWaterson2005} shows a rival's existing
presence predicts a firm's own entry dynamically; \citet{BerryReiss2007} surveys the tradition. None of this
literature has been applied to cannabis border markets specifically, and none pairs an entry-regime moderator
with an exogenous demand shock of the kind cross-border legalization timing provides.

The gap is specific and modest, not a new mechanism. The cross-border demand literature has established that
demand moves across a cannabis-policy border; it has not asked where that demand goes once it arrives. The
entry-regulation literature has established that caps distort market structure and rents; it has not had an
exogenous demand shock to test that distortion against. This paper bridges those two literatures with a
design built for that purpose, and reports, rather than overstates, what the resulting test does and does not
show.

\section{Institutional Background}
\label{sec:institutional}

\input{tables/table_state_summary}

Table~\ref{tab:state_summary} summarizes each sample state's legalization timing, border exposure, data
sources, and entry-regulation regime.

\textbf{Colorado} imposes no state-level numerical cap on retail marijuana licenses -- Amendment 64 (2012)
delegates licensing authority to local governments entirely. But \citet{PayanBrownSong2021}'s county-level
census of Colorado and Washington's local marijuana ordinances finds that, as of April 2019, 48.4 percent of
Colorado counties banned all marijuana facilities outright, and only 23.4 percent allowed every facility type
without restriction. We classify Colorado as capped for its entire panel on this basis: a state law that says
``no cap'' does not describe a free-entry market if local governments can and routinely do close most of the
state's geography to retail entirely.

\textbf{Washington} caps entry directly. Initiative 502 (2012) set an initial statewide ceiling of 334 retail
licenses; the Washington State Liquor and Cannabis Board raised it to 556 by emergency rule on January 6, 2016,
and a social-equity program began issuing up to 52 additional licenses in 2024. The cap level changes three
times over the sample window, but the state is never free-entry at any point in it.

\textbf{Oregon} is the one state whose regime genuinely switches mid-panel, and we treat it separately
(Section~\ref{sec:results}). Retail licensing was open from the market's October 2015 launch until a producer
moratorium (2018--2019) and a full application freeze effective January 2, 2022 (House Bill 4016, March 2022).
House Bill 4121 (March 2024) replaced the freeze with a permanent formula, one retail license per 7,500
residents aged 21 and over. We code Oregon free through 2021 and capped from 2022 onward.

\textbf{Montana}'s enabling statute (HB 701, retail effective January 1, 2022) imposes no numerical ceiling on
dispensary locations, but restricts new licenses to conversions of existing medical marijuana providers and
requires counties that voted against the state's Initiative 190 (November 2020) to affirmatively opt in before
a dispensary can locate there -- binding from the market's first day, not merely from a later moratorium
(Senate Bill 27, May 2025). Montana is coded capped for its entire panel on the opt-in requirement alone.

\textbf{New Mexico} is the cleanest free-entry state in the sample. The Cannabis Regulation Act (NMSA 1978
\S 26-2C-6(I)) explicitly bars the state's Cannabis Control Division from limiting license counts, and
licensing is administrative -- a 90-day issue-or-deny window, no lottery or merit-scoring step. Local
governments retain zoning authority but cannot ban cannabis retailers outright. New Mexico is coded free
throughout.

\section{Data}
\label{sec:data}

\subsection{Revenue and dispensary-count data}

Colorado's monthly county-level marijuana sales come from the Colorado Department of Revenue's Socrata series
(2014--2025). Washington's come from the Liquor and Cannabis Board's fiscal-year county sales and excise-tax
reports (FY2015--FY2025). Oregon's monthly county sales come from the Oregon Liquor and Cannabis Commission's
own Socrata series (2017--2025); Oregon's revenue-per-store outcome exists only as a 2025 cross-section, since
OLCC's license registry records a license's most recent renewal date rather than its original opening date.
Montana's monthly county sales were scraped from the Department of Revenue's static report pages, no bulk
download existing (2022--2025); Montana's revenue-per-store outcome has the same cross-section-only limitation
as Oregon's. New Mexico's data come from the Cannabis Control Division's Cannabis Reporting Online Portal,
which publishes monthly sales both by city and, uniquely among the five states, by individual licensee, so
New Mexico's revenue-per-store outcome is a genuine historical panel (2022--2026), not a cross-section.
Dispensary counts for Colorado and Washington come from a geocoded license dataset, point-in-polygon assigned
to counties.

\subsection{Time-varying border classification}
\label{sec:timevarying}

Prior versions of this analysis, including our own, classified border exposure using a single legal-status
snapshot at a fixed reference date -- incorrect for any county-year in which a neighboring state's own status
differed from that snapshot. Three cases in this sample are large enough to matter. Colorado borders New
Mexico (legal 2022-04-01); for Colorado county-years before that date, eight Colorado counties (Costilla,
Archuleta, Conejos, Las Animas, Montezuma, La Plata, Rio Grande, and Alamosa) fall within 70 miles of a New
Mexico county, several not within 70 miles of any other nonlegal state. Washington borders Oregon (legal
2015-10-01); fourteen Washington counties, including populous Clark County across the Columbia River from
Portland, fall within 70 miles of Oregon, bearing on Washington's earliest fiscal-year observations. Oregon
itself borders California (legal 2018-01-01) and Nevada (legal 2017-07-01); three Oregon counties (Josephine,
Curry, Jackson) fall within 70 miles of California, bearing on Oregon's first panel year, 2017.

We recompute distance to the nearest nonlegal-state county separately for every observation year, using each
neighbor's own retail-legalization date. This reclassifies 41 of 1,314 per-capita county-year observations and
38 of 995 revenue-per-store observations relative to the static alternative, concentrated in the early years
of Washington's and Colorado's panels and Oregon's first year.

\subsection{Entry-regime coding}

We code a binary \texttt{Capped} indicator at the state-year level: Colorado, Washington, and Montana are
capped throughout; Oregon is free through 2021 and capped from 2022; New Mexico is free throughout. Limitations
of this state-year measure are discussed in Section~\ref{sec:discussion}.

\subsection{Local- versus state-level entry barriers}
\label{sec:localbarrier}

We also construct an alternative classification, \texttt{LocalBarrier}, isolating specifically \emph{which
level of government} imposes the binding entry restriction, holding fixed the same institutional facts used to
build \texttt{Capped}. Colorado and Montana are coded 1 (the restriction is county-level -- opt-out bans in
Colorado, opt-in requirements in Montana); Washington, Oregon, and New Mexico are coded 0 (Washington's cap is
set directly by state statute with no local mechanism; Oregon's freeze is a state-level action; New Mexico has
no restriction). This reclassifies Washington relative to \texttt{Capped} while leaving Colorado and Montana
unchanged, and is used as a robustness check on \texttt{Capped} in Section~\ref{sec:results}.

\subsection{Dispensary-level panel: New Mexico}
\label{sec:nmdispensary}

Every result above uses county-year aggregates. To test whether the county-level border premium survives at
the level of the individual seller, we separately construct a genuine dispensary-level panel for New Mexico --
the only sample state where individual-firm revenue is available historically. New Mexico's Cannabis Control
Division publishes 51 monthly licensee-level files (April 2022 onward), each listing every active licensee's
street address and reported sales. We identify a dispensary by its normalized street address rather than its
licensee name, since the same location can appear under different names across months (ownership changes,
rebrands), while different addresses under the same name are genuinely different stores. This yields 1{,}253
unique dispensary locations. Each dispensary's entry month is the first calendar month its address appears;
231 dispensaries (18.4 percent) are already present in the panel's first month (April 2022, the adult-use
launch date) and are flagged as left-censored (almost certainly pre-existing medical-program dispensaries,
whose true entry date is unobserved), while the remaining 1{,}022 are true new entrants whose opening month we
observe directly. County assignment reuses the same ZIP-to-county crosswalk already validated for the
county-level New Mexico panel; we do not build a separate address-level geocoding pipeline, since no
address-level coordinate source exists for New Mexico in our data (unlike Colorado and Washington, where a
separately assembled geocoded license file gives coordinates and opening dates but not store-level revenue).
The regression sample drops a small number of dispensary-years with zero reported revenue, leaving 1{,}249
dispensaries and 3{,}703 dispensary-year observations. Because New Mexico's border exposure and entry regime
are both time-invariant across the sample, this exercise cannot test the entry-regulation interaction; it
serves as a firm-level robustness check on the existence of the border premium itself.

\section{Empirical Strategy}
\label{sec:methods}

The main specification is
\[
y_{ist} = \alpha + \beta_1 \, \text{Border}_{ist} + \beta_2 \, \text{Capped}_{st} + \beta_3\,
\text{Border}_{ist} \times \text{Capped}_{st} + \gamma \ln(\text{population}_{ist}) + \delta_t + \eta_i + \varepsilon_{ist}
\]
estimated separately for log(revenue per dispensary store), log(revenue per capita), and log(dispensary count
per 1{,}000 residents). $\beta_3$ is the coefficient of interest, testing whether the border premium differs
by entry regime. State fixed effects $\eta_i$ are included for revenue per capita, where Oregon's within-state
regime switch identifies $\beta_3$ independent of the cross-state contrast, but omitted for revenue per store
and density: Oregon's per-store data is a single 2025 cross-section, so \texttt{Capped} has no within-state
variation for any state in that sample and becomes collinear with state fixed effects, which we verified
directly produces an unstable, uninterpretable coefficient. Standard errors are clustered by county throughout
(165--178 clusters), with HC1 standard errors reported alongside for comparison. Every specification is also
estimated separately by state, and Oregon's own free/capped switch is tested a third way, splitting its own
panel at the January 2022 freeze.

\section{Results}
\label{sec:results}

\input{tables/table_main_pooled}

Table~\ref{tab:main_pooled} reports the main pooled specification. The border premium is positive and
significant on all three outcomes: revenue per store is 48 percent higher in border counties (county-clustered
$p=0.11$), revenue per capita is 193 percent higher ($p<0.001$), and dispensary density is 101 percent higher
($p<0.001$). The interaction with capped entry is a 54\% relative reduction in the border premium's
multiplier on revenue per capita (from roughly 193\% to roughly 34\% in level terms, $p=0.009$) and a 56\%
relative reduction on dispensary density ($p<0.001$), but a precise null on revenue per store (3.5 percent,
$p=0.90$). Figure~\ref{fig:premium_by_regime} shows the pattern directly.

\begin{figure}[htbp]
\centering
\includegraphics[width=0.85\textwidth]{figures/fig1_premium_by_regime.png}
\caption{Border premium by outcome and entry regime, with 95\% confidence intervals.}
\label{fig:premium_by_regime}
\end{figure}

\input{tables/table_state_specific}

Table~\ref{tab:state_specific} reports the same outcomes estimated separately within each state.
Colorado's own density coefficient is negative and significant ($-27$ percent) despite its large positive
revenue-per-store premium ($+81$ percent), matching the pooled interaction's prediction at the state level.
New Mexico's density coefficient is strongly positive ($+109$ percent). Montana's and Oregon's estimates are
noisier, reflecting small per-store and density subsamples (23 and 30 observations respectively).

\begin{figure}[htbp]
\centering
\includegraphics[width=0.85\textwidth]{figures/fig2_state_forest.png}
\caption{State-specific border coefficients by outcome, with 95\% confidence intervals (HC1).}
\label{fig:state_forest}
\end{figure}

\input{tables/table_oregon_beforeafter}

Table~\ref{tab:oregon_beforeafter} and Figure~\ref{fig:oregon_beforeafter} split Oregon's own revenue-per-capita
panel at its January 2022 freeze. The border premium is larger in Oregon's free period (+173 percent,
$p=0.007$) than its capped period (+106 percent, $p=0.10$), the same direction as the pooled interaction,
though the formal within-Oregon interaction is not itself significant ($p=0.50$ HC1), which we attribute to
Oregon having only three border counties rather than to a contradiction of the pooled result.

\begin{figure}[htbp]
\centering
\includegraphics[width=0.7\textwidth]{figures/fig3_oregon_beforeafter.png}
\caption{Oregon's revenue-per-capita border premium, before/after the January 2022 freeze.}
\label{fig:oregon_beforeafter}
\end{figure}

\begin{figure}[htbp]
\centering
\includegraphics[width=0.72\textwidth]{figures/fig5_distance_gradient.png}
\caption{Revenue per store against log-distance to the nearest nonlegal-state county, residualized.}
\label{fig:distance_gradient}
\end{figure}

Replacing the binary border indicator with continuous log-distance recovers the same pattern (slope $-0.265$,
$p<0.001$; Figure~\ref{fig:distance_gradient}).

\begin{figure}[htbp]
\centering
\includegraphics[width=0.72\textwidth]{figures/fig6_captive_population.png}
\caption{Revenue per store against log nonlegal-side population within 70 miles, residualized.}
\label{fig:captive_population}
\end{figure}

The total population of nonlegal-state counties within 70 miles, the most robust mechanism evidence in our
earlier two-state analysis, weakens to a precise null pooling all five states (slope $+0.015$, HC1 $p=0.54$,
county-clustered $p=0.77$; Figure~\ref{fig:captive_population}) -- reported directly rather than retaining the
more favorable two-state framing.

\subsection{Robustness: local- versus state-level entry barriers, and a dispensary-level check}

\input{tables/table_localbarrier_pooled}

Table~\ref{tab:localbarrier_pooled} re-estimates the border-by-regime interaction using \texttt{LocalBarrier}
(Section~\ref{sec:localbarrier}) in place of \texttt{Capped}. The density interaction survives and remains
significant under county clustering ($-41.5$ percent, $p=0.037$); the revenue-per-capita and revenue-per-store
interactions are significant under HC1 but lose significance once clustered by county (per capita: $p=0.35$;
per store: $p=0.17$) -- a genuine caveat on the main classification, reported directly. Figures
\ref{fig:localbarrier_percapita} and \ref{fig:localbarrier_perstore} show a pooled year-by-year event study of
the same comparison (county population, log(1+captive population), state and year fixed effects, clustered by
county, Oregon excluded): on revenue per capita, local-barrier states (CO, MT) show a border premium roughly
100 percentage points above non-local-barrier states (WA, NM) through 2019, converging by 2024; on revenue per
store, the two groups track closely throughout.

\begin{figure}[htbp]
\centering
\includegraphics[width=0.85\textwidth]{figures/fig11_localbarrier_v2_percapita.png}
\caption{Border premium on revenue per capita, by year, local-barrier vs.\ non-local-barrier states.}
\label{fig:localbarrier_percapita}
\end{figure}

\begin{figure}[htbp]
\centering
\includegraphics[width=0.85\textwidth]{figures/fig12_localbarrier_v2_perstore.png}
\caption{Border premium on revenue per store, by year, local-barrier vs.\ non-local-barrier states.}
\label{fig:localbarrier_perstore}
\end{figure}

\input{tables/table_nm_dispensary_eventstudy}

Table~\ref{tab:nm_dispensary_eventstudy} and Figure~\ref{fig:nm_dispensary} re-estimate the border premium
directly on New Mexico's individual dispensaries (Section~\ref{sec:nmdispensary}) rather than on county
totals. The firm-level premium is large, positive, and significant in every year (113 to 184 percent, all
$p<0.02$); a pooled single-coefficient version recovers $+141.4$ percent ($p=0.001$), close to New Mexico's
county-level contribution to the main estimate. Border dispensaries are also disproportionately new entrants
(87.2 percent vs.\ 78.3 percent interior), though descriptively rather than causally, since New Mexico's own
regime never varies. This exercise cannot test the entry-regulation interaction, but confirms the border
premium is not a county-aggregation artifact in the one state where individual-firm revenue is observable.

\begin{figure}[htbp]
\centering
\includegraphics[width=0.75\textwidth]{figures/fig13_nm_dispensary_eventstudy.png}
\caption{Border premium on individual-dispensary revenue, by year, New Mexico firm-level panel
($n=1{,}249$ dispensaries).}
\label{fig:nm_dispensary}
\end{figure}

\subsection{Economic magnitude: an illustrative revenue-gap accounting}
\label{sec:magnitude}

The interaction estimate implies a large difference in level terms, not just in the log coefficients
themselves: converting $\beta_1$ and $\beta_1+\beta_3$ separately, the free-entry border premium is
$+193.3$ percent while the capped-entry border premium is $+34.3$ percent, so a capped-state border county
would need $2.18\times$ its actual revenue to match what the free-entry premium implies for a county of its
size. Applying this ratio to actual 2025 border-county revenue in the three capped states in our sample
(Colorado, Washington, Montana; Table~\ref{tab:magnitude}) implies roughly \$480 million in aggregate revenue,
and roughly \$117 million in state tax revenue at each state's own statutory cannabis tax rate (Colorado 15
percent, Washington 37 percent, Montana 20 percent), that these states' border counties would need to gain to
match the free-entry premium's implied level, on a single year's data.

\begin{table}[htbp]
\centering
\caption{Illustrative revenue-gap accounting, capped-state border counties, 2025}
\label{tab:magnitude}
\small
\begin{tabular}{lccccc}
\toprule
State & Actual revenue & Counterfactual & Revenue gap & Tax rate & Tax gap \\
\midrule
Colorado & \$195.9M & \$427.8M & \$231.9M & 15\% & \$34.8M \\
Washington & \$161.6M & \$352.8M & \$191.2M & 37\% & \$70.8M \\
Montana & \$48.3M & \$105.6M & \$57.2M & 20\% & \$11.4M \\
\midrule
Total & \$405.8M & \$886.2M & \$480.4M & -- & \$117.0M \\
\bottomrule
\end{tabular}
\begin{minipage}{0.9\textwidth}
\vspace{0.3em}
\footnotesize
\textit{Notes:} Counterfactual = actual $\times$ 2.1837, the ratio of $(1+\text{free premium})$ to
$(1+\text{capped premium})$ implied by the main pooled interaction (Table~\ref{tab:main_pooled}). Montana's
tax rate is its state adult-use cannabis excise tax; local-option taxes (up to an additional 3\% in some
Montana counties) and Colorado's separate wholesale excise tax are excluded.
\end{minipage}
\end{table}

We present this explicitly as an illustrative accounting exercise, not a causal welfare estimate, for three
reasons. First, it assumes external validity: that Colorado, Washington, and Montana's border counties would
earn revenue at New Mexico's and free-period Oregon's own rate if their entry regime changed, an extrapolation
the cross-state design does not directly identify, and Oregon's own within-state before/after test
(Section~\ref{sec:results}) points the same direction but is not independently significant. Second, the figure
is a revenue gap, not a welfare loss: it says nothing about whether the difference reflects genuine unmet
demand (a real economic loss, e.g.\ continued cross-border leakage) or a different distribution of similar
underlying demand, since higher revenue per store under a binding cap is also consistent with the market
power documented by \citet{HollenbeckUetake2021} for Washington's own licensing regime, and our data contain
no prices or quantities that would let us adjudicate between these. Third, it is a single-year (2025) snapshot,
not a cumulative or discounted figure. We report it as an order of magnitude for a policy audience, not as a
precise estimate.

\section{Discussion}
\label{sec:discussion}

The defensible claim here is narrower than the paper's opening hypothesis. Cross-border cannabis demand
produces a measurable premium in border-county markets on every outcome we observe. Entry regulation
moderates that premium on the quantity margin specifically -- store count and per-capita revenue -- not on
price at any single store. The original hypothesis predicted regulation would redirect the premium from
quantity into price; the data show regulation blocks the quantity response without redirecting anything, and
the price-margin premium exists regardless of regime.

Three limitations. First, the entry-regime variable is state-year, not county-year, though the institutional
detail motivating Colorado's and Montana's capped classification is itself local; a county-level restriction
measure is a natural extension we did not build here. Second, Oregon's own before/after test is directionally
consistent with the pooled result but not independently significant, and should be weighted as corroborating
rather than confirmatory. Third, the captive-population result reverses statistical significance once the
sample expands from two states to five states, a reminder that small-sample mechanism evidence does not
automatically generalize. Fourth, an alternative cut of the same entry-regime classification -- isolating
specifically local-government-imposed restrictions (Section~\ref{sec:localbarrier}) -- preserves the density
result under county clustering but not the revenue-per-capita or revenue-per-store interactions
(Table~\ref{tab:localbarrier_pooled}), a reminder that the interaction's precise magnitude, though not its
direction on the quantity margin, is sensitive to how the capped/free distinction is drawn. Set against these
limitations, the New Mexico dispensary-level panel (Section~\ref{sec:nmdispensary}) is a genuine strength: it
confirms the border premium itself is not an artifact of county-level aggregation.

\section{Conclusion}
\label{sec:conclusion}

Uneven legalization creates real spatial demand spillovers in legal cannabis markets. Entry regulation --
coded from actual statutes and local-government authority, not from a state's own description of itself --
determines whether that premium flows into market expansion: capped-entry states see a substantially smaller
premium on revenue per capita and dispensary density, while free-entry states' extra demand shows up as more
stores. The premium on revenue at any single store does not depend on regulatory regime, a real departure from
the paper's starting hypothesis that we report rather than obscure. The result connects two literatures --
cross-border cannabis demand and cannabis entry regulation -- that have so far developed independently.

\bibliographystyle{apalike}
\bibliography{references}

\end{document}
